\section{Линеаризация обратной связью}
\label{sec:linearization}

\subsection{Идея метода: компенсировать нелинейность управлением}

\textit{L4\_1, слайды 1--9; L4\_2, слайды 1--2.}
Линеаризация обратной связью использует измеренное состояние, чтобы
скомпенсировать нелинейные слагаемые и получить заданную линейную
динамику. Компенсация выполняется в самом законе управления.
Например, для маятника, записанного относительно требуемого положения,
\[
 \dot x_1=x_2,\qquad
 \dot x_2=-a[\sin(x_1+\delta)-\sin\delta]-bx_2+cu,
 \qquad c\ne0,
\]
выбор
\[
 u=\frac{a}{c}[\sin(x_1+\delta)-\sin\delta]+\frac{v}{c}
\]
даёт точно $\dot x_2=-bx_2+v$. Затем выбираем
$v=-k_1x_1-k_2x_2$ и получаем
\[
 \dot{\bm x}=\begin{pmatrix}0&1\\-k_1&-(b+k_2)\end{pmatrix}\bm x,
 \qquad
 p(\lambda)=\lambda^2+(b+k_2)\lambda+k_1.
\]
Для гурвицевости достаточно и необходимо $k_1>0$, $b+k_2>0$.
Нелинейная часть компенсирует гравитацию, а линейная часть формирует
возвращающий момент и демпфирование.

\textbf{Дополнение для защиты.} При обычной линеаризации по Тейлору
нелинейную модель приближённо заменяют линейной около равновесия.
Здесь при точной модели получается точное равенство на всей области,
где определены обратная связь и замена координат. Поэтому нужно
проверять не малость отклонений, а обратимость преобразования,
ненулевой коэффициент при управлении и поведение внутренних состояний.
Измерения состояния, достаточная точность модели и возможность
реализовать рассчитанное управление являются условиями применения.

\subsection{Полная линеаризация и выбор матрицы усиления}

\textit{L4\_1, слайды 10--18.}
Для системы, аффинной по управлению,
\[
 \dot{\bm x}=f(\bm x)+G(\bm x)\bm u,
 \qquad f(0)=0,
\]
предполагается существование замены $\bm z=T(\bm x)$, приводящей
систему к управляемой форме
\[
 \dot{\bm z}=A\bm z+B\gamma(\bm x)[\bm u-\alpha(\bm x)].
\]
Здесь $(A,B)$ --- управляемая пара, а квадратная матрица
$\gamma(\bm x)$ невырождена в рассматриваемой области.
\emph{Скобки существенны}: $\gamma$ умножает всю разность
$[\bm u-\alpha]$. Тогда
\[
 \bm u=\alpha(\bm x)+\gamma(\bm x)^{-1}\bm v,
 \qquad \bm v=-K\bm z
\]
даёт $\dot{\bm z}=(A-BK)\bm z$. В исходных координатах
\[
 \boxed{\bm u=\alpha(\bm x)-\gamma(\bm x)^{-1}KT(\bm x)},
 \qquad
 \dot{\bm x}=f(\bm x)+G(\bm x)
 [\alpha(\bm x)-\gamma(\bm x)^{-1}KT(\bm x)].
\]
В этой формуле $\bm z$ обозначает \emph{весь} преобразованный вектор,
как в L4\_1. Ниже внешние координаты нормальной формы обозначаются
$\bm\xi$, как в L4\_2. Вспомогательный вход на части слайдов L4\_1
обозначен $\nu$; далее используется $v$ из L4\_2 и задания.

Управляемость позволяет выбрать $K$ по желаемым собственным значениям
$A-BK$. Для цепочки $\rho$ интеграторов
$K=(k_1,\ldots,k_\rho)$ даёт полином
\[
 p(\lambda)=\lambda^\rho+k_\rho\lambda^{\rho-1}+\cdots+k_2\lambda+k_1.
\]
То есть коэффициенты $K$ выписываются по возрастанию степеней
полинома, начиная со свободного члена. Во внешней динамике нормальной
формы длина цепочки равна относительной степени $\rho$.

\textbf{Дополнение для защиты.} Невырожденный якобиан обеспечивает
лишь локальную обратимость. Для глобального результата отдельно
нужны глобальное обратное преобразование и сохранение рассматриваемой
области траекториями. Глобальный диффеоморфизм, переводящий нуль в нуль,
переносит глобальную асимптотическую устойчивость. Глобальная
экспоненциальная оценка в исходной евклидовой норме дополнительно
требует равномерной сравнимости норм $\norm{T(\bm x)}$ и
$\norm{\bm x}$; одна гладкая обратимость этого не гарантирует.

\subsection{Производные Ли и относительная степень}

\textit{L4\_2, слайды 2--6.}
Для одного входа и одного выхода рассматривается
\[
 \dot{\bm x}=f(\bm x)+g(\bm x)u,\qquad y=h(\bm x),
 \qquad \bm x\in D\subset\mathbb R^n.
\]
Функция $h$ и векторные поля $f,g$ должны быть достаточно гладкими.
Производная Ли функции вдоль поля --- её скорость изменения при
движении вдоль этого поля:
\[
 L_fh=\frac{\partial h}{\partial\bm x}f,
 \qquad L_gh=\frac{\partial h}{\partial\bm x}g,
 \qquad \dot y=L_fh+L_gh\,u.
\]
Градиент здесь записан строкой, поэтому произведение имеет скалярный
результат. Повторные производные определяются рекурсивно:
\[
 L_f^0h=h,\qquad L_f^kh=L_f(L_f^{k-1}h),\qquad
 L_gL_fh=\frac{\partial(L_fh)}{\partial\bm x}g.
\]
Если $L_gh=0$ во всей рассматриваемой области, то $\dot y=L_fh$,
и можно ещё раз дифференцировать:
\[
 \ddot y=L_f^2h+L_gL_fh\,u.
\]
Продолжаем, пока управление впервые не появится с ненулевым
коэффициентом. Система имеет относительную степень $\rho$ на
$D_0\subset D$, если
\[
 L_gL_f^{i-1}h=0\quad (i=1,\ldots,\rho-1),\qquad
 L_gL_f^{\rho-1}h\ne0\quad \text{для всех }\bm x\in D_0.
\]
Тогда $1\le\rho\le n$ и
\[
 y^{(\rho)}=L_f^\rho h+L_gL_f^{\rho-1}h\,u,
 \qquad
 \boxed{u=\frac{-L_f^\rho h+v}{L_gL_f^{\rho-1}h}}
 \quad\Longrightarrow\quad y^{(\rho)}=v.
\]
Относительная степень --- число дифференцирований \emph{выбранного
выхода} до появления входа, а не число состояний вообще.
Буква $\rho$ здесь не является $\varrho(\bm x)$ --- оценкой
неопределённости из лекции о скользящем режиме.

\textbf{Примеры со слайда 5.} Для уравнения Ван-дер-Поля
\[
 \dot x_1=x_2,\qquad
 \dot x_2=-x_1+\epsilon(1-x_1^2)x_2+u
\]
получаем три разных результата: при $y=x_1$ имеем $\rho=2$;
при $y=x_2$ имеем $\rho=1$; при $y=x_1+x_2^2$ коэффициент
$L_gh=2x_2$, поэтому $\rho=1$ только на $x_2\ne0$.
Для системы $\dot x_1=x_1$, $\dot x_2=x_2+u$, $y=x_1$
все производные выхода равны $y$: управление на выход не влияет,
и относительная степень не определена.

Третий пример лекции:
\[
 \dot x_1=-ax_1+u,\quad
 \dot x_2=-bx_2+k-cx_1x_3,\quad
 \dot x_3=\theta x_1x_2,\quad y=x_3.
\]
Здесь $L_gh=0$, $L_fh=\theta x_1x_2$ и
$L_gL_fh=\theta x_2$. При $\theta\ne0$ относительная степень
равна двум на $x_2\ne0$. Пример показывает, почему область
определения нужно указывать вместе с числом~$\rho$.

\textbf{Дополнение для защиты.} Проверка только в одной точке
недостаточна: предыдущие коэффициенты должны тождественно обращаться
в нуль в окрестности. Если последний коэффициент занулился на
гиперповерхности, нельзя автоматически объявить там относительную
степень на единицу большей. После фактического появления $u$
дальнейшее дифференцирование может породить $\dot u$ и нелинейные
функции управления.

\subsection{Нормальная форма, внутренние координаты и диффеоморфизм}

\textit{L4\_2, слайды 7--8.}
Внешние координаты естественно выбрать как выход и его производные,
которые ещё не содержат управления:
\[
 \bm\xi=\psi(\bm x)=
 (h,L_fh,\ldots,L_f^{\rho-1}h)^T.
\]
Они имеют размерность $\rho$. Если $\rho<n$, до полного вектора
не хватает $n-\rho$ внутренних координат
$\bm\eta=\varphi(\bm x)$. Их ищут из условий
\[
 L_g\varphi_i=\frac{\partial\varphi_i}{\partial\bm x}g=0,
 \qquad
 T(\bm x)=\begin{pmatrix}\varphi(\bm x)\\\psi(\bm x)\end{pmatrix}
 \text{ --- диффеоморфизм}.
\]
Первое условие убирает непосредственное управление из $\dot\eta_i$:
\[
 \dot{\bm\eta}=f_0(\bm\eta,\bm\xi)
 =\left.\frac{\partial\varphi}{\partial\bm x}f(\bm x)
 \right|_{\bm x=T^{-1}(\bm\eta,\bm\xi)}.
\]
Зависимость от внешних координат при этом обычно сохраняется.
Для одного входа подходящие внутренние координаты локально
существуют при постоянной относительной степени. Если $\rho=n$,
внутренних координат нет, а $T=\psi$ локально обратимо.

Обозначим
\[
 \gamma(\bm x)=L_gL_f^{\rho-1}h(\bm x),\qquad
 \alpha(\bm x)=-\frac{L_f^\rho h(\bm x)}{\gamma(\bm x)}.
\]
Тогда нормальная форма \emph{исходной} системы имеет вид
\[
 \dot{\bm\eta}=f_0(\bm\eta,\bm\xi),\qquad
 \dot{\bm\xi}=A_c\bm\xi+B_c\gamma(\bm x)[u-\alpha(\bm x)],
 \qquad y=C_c\bm\xi,
\]
где $A_c$ имеет единицы над главной диагональю и нули в остальных
элементах, $B_c=(0,\ldots,0,1)^T$, $C_c=(1,0,\ldots,0)$.
Последняя строка равна $\dot\xi_\rho=L_f^\rho h+\gamma u$,
что позволяет проверить знак $\alpha$.
Обратная связь $u=\alpha+\gamma^{-1}v$ даёт
\[
 \dot{\bm\eta}=f_0(\bm\eta,\bm\xi),\qquad
 \dot{\bm\xi}=A_c\bm\xi+B_cv.
\]
Таким образом, замена координат и замена входа --- два отдельных
шага. Линеаризация по входу-выходу при $\rho<n$ линеаризует только
внешнюю цепочку.

\textbf{Дополнение для защиты.} Диффеоморфизм --- взаимно однозначное
гладкое преобразование с гладким обратным. По теореме об обратной
функции условие $\det DT(\bm x_0)\ne0$ гарантирует его существование
в некоторой окрестности $\bm x_0$. Оно не гарантирует однозначность
на большой области: например, квадрат $x_3^2$ теряет знак $x_3$.
При проверке глобальности полезнее явно выписать $T^{-1}$ и его область.

\subsection{Нуль-динамика и минимально-фазовость}

\textit{L4\_2, слайды 9--11.}
Нуль-динамика --- внутреннее движение, совместимое с тождественно
нулевым выходом. Для этого нужны $\bm\xi\equiv0$ и $v\equiv0$,
поэтому
\[
 \dot{\bm\eta}=f_0(\bm\eta,0).
\]
В исходных координатах оно происходит на множестве
\[
 Z^*=\{\bm x\in D_0:h=L_fh=\cdots=L_f^{\rho-1}h=0\},
 \qquad u^*(\bm x)=\left.\alpha(\bm x)\right|_{Z^*},
\]
и определяется ограничением поля
$f+g\alpha$ на $Z^*$. Нулевой выход не означает нулевое состояние:
на $Z^*$ могут оставаться свободные внутренние переменные.

Если $T$ переводит равновесие в $(\bm\eta,\bm\xi)=(0,0)$,
система называется \emph{минимально-фазовой}, когда нуль внутренней
системы асимптотически устойчив. Обычной устойчивости недостаточно.
Это хорошо видно в примере Ван-дер-Поля с $y=x_2$:
\[
 \eta=x_1,\quad\xi=x_2,\qquad
 \dot\eta=\xi,\quad
 \dot\xi=-\eta+\epsilon(1-\eta^2)\xi+u.
\]
При $\xi=0$ и сохраняющем нулевой выход управлении $u=\eta$
получаем $\dot\eta=0$. Любое внутреннее состояние сохраняется,
так что система не минимально-фазовая.

\textbf{Дополнение для защиты.} Даже минимально-фазовость сама по
себе является локальным свойством нуль-динамики и не заменяет
проверку всей системы при $\bm\xi\ne0$. Для глобальной стабилизации
нужно исключить уход из области координат и неограниченный рост
внутренних состояний. В лабораторной задаче 2 это сделано точным
приведением \emph{всех} новых координат к линейной гурвицевой системе.

\subsection{Последний пример L4\_2: полезный результат и опечатка}

\textit{L4\_2, слайд 12; формулы проверены по изображению слайда.}
На слайде напечатано
\[
 \dot x_1=-x_1+\frac{2+x_2^3}{1+x_2^3}u,\qquad
 \dot x_2=x_3,\qquad\dot x_3=x_1x_3+u,\qquad y=x_2.
\]
Отсюда $\dot y=x_3$, $\ddot y=x_1x_3+u$, поэтому
$\rho=2$, $\gamma=1$, $\alpha=-x_1x_3$.
На $Z^*=\{x_2=x_3=0\}$ имеем $u^*=0$ и
$\dot x_1=-x_1$. Следовательно, в окрестности начала система
минимально-фазовая.

\textbf{Дополнение для защиты: проверка согласованности слайда.}
У напечатанного поля знаменатель обращается в нуль при $x_2=-1$:
утверждение о гладкой модели во всём $\mathbb R^3$ здесь неверно.
Кроме того, предложенная там функция
$\varphi=-x_1+x_3+\arctan x_3$ даёт
\[
 L_g\varphi=-\frac{2+x_2^3}{1+x_2^3}
              +1+\frac{1}{1+x_3^2}\not\equiv0.
\]
Для \emph{напечатанной} системы корректно, например, взять
\[
 a(x_2)=\frac{2+x_2^3}{1+x_2^3},\qquad
 \eta=x_1-a(x_2)x_3,\quad \xi_1=x_2,\quad\xi_2=x_3.
\]
Здесь $L_g\eta=0$, $\det DT=1$, а обратное преобразование
$x_1=\eta+a(\xi_1)\xi_2$ определено при $\xi_1\ne-1$.

Вероятное исправление исходного примера --- коэффициент
$(2+x_3^2)/(1+x_3^2)$ перед $u$. Тогда именно функция со слайда
\[
 \eta=-x_1+x_3+\arctan x_3
\]
удовлетворяет $L_g\eta=0$; преобразование
$(\eta,\xi_1,\xi_2)=(\eta,x_2,x_3)$ имеет определитель $-1$
и глобальное обратное $x_1=-\eta+\xi_2+\arctan\xi_2$.
На нулевом выходе $\eta=-x_1$ и $\dot\eta=-\eta$.
Этот согласованный вариант объясняет замысел примера, но не
подменяет собой напечатанное условие.

\subsection{Ошибки модели и запас робастности}

\textit{L4\_1, слайды 19--22.}
Точная компенсация требует истинных $\alpha,\gamma,T$.
На практике используются номинальные модели с крышками:
\[
 \bm u=\widehat\alpha(\bm x)
       -\widehat\gamma(\bm x)^{-1}K\widehat T(\bm x).
\]
В истинных координатах $\bm z=T(\bm x)$ ошибка остаётся как
добавка во входном канале:
\[
 \dot{\bm z}=(A-BK)\bm z+B\delta(\bm z),
\]
\[
 \delta(T(\bm x))=\gamma(\bm x)
 \left[\widehat\alpha(\bm x)-\alpha(\bm x)
 +\gamma(\bm x)^{-1}KT(\bm x)
 -\widehat\gamma(\bm x)^{-1}K\widehat T(\bm x)\right].
\]
Сначала группируем ошибку компенсации и ошибку линейной обратной
связи внутри скобок, затем умножаем всё на истинную $\gamma$.
Равенство понимается после подстановки $\bm x=T^{-1}(\bm z)$.

Пусть $A_K=A-BK$ гурвицева. Решаем уравнение Ляпунова
\[
 PA_K+A_K^TP=-I,\qquad P=P^T>0,
\]
и берём $V=\bm z^TP\bm z$. Тогда
\[
 \dot V=-\norm{\bm z}^2+2\bm z^TPB\delta(\bm z)
 \le-\norm{\bm z}^2+2\norm{PB}\norm{\bm z}\norm{\delta(\bm z)}.
\]
Если на всей области $\norm{\delta(\bm z)}\le k\norm{\bm z}$ и
\[
 0\le k<\frac{1}{2\norm{PB}},
\]
то $\dot V\le-[1-2\norm{PB}k]\norm{\bm z}^2<0$ при
$\bm z\ne0$. При глобальном выполнении условий получаем глобальную
экспоненциальную устойчивость в $\bm z$. Здесь $k$ --- оценка ошибки
модели, а $K$ --- матрица регулятора: их нельзя смешивать.

\textbf{Дополнение для защиты.} Условие $\delta(0)=0$ уже следует из
оценки $\norm\delta\le k\norm{\bm z}$. Постоянная ошибка компенсации
в нуле этой оценке не удовлетворяет. При
$\norm\delta\le k\norm{\bm z}+\varepsilon$, $\varepsilon>0$,
получается
\[
 \dot V\le-q\norm{\bm z}^2+2\norm{PB}\varepsilon\norm{\bm z},
 \qquad q=1-2\norm{PB}k>0.
\]
Производная отрицательна вне шара радиуса
$2\norm{PB}\varepsilon/q$. Это даёт предельную ограниченность
в некоторой окрестности; точную сходимость к исходному нулю
такая оценка не доказывает.

\subsection{Маятник: точная и номинальная компенсация}

\textit{L4\_1, слайды 23--29.}
Физическая модель и требуемое положение:
\[
 \ddot\theta=-a\sin\theta-b\dot\theta+cT,\qquad \theta=\delta.
\]
Для равновесия требуется $T_{ss}=(a/c)\sin\delta$.
После замены $x_1=\theta-\delta$, $x_2=\dot\theta$,
$u=T-T_{ss}$ получается система из начала главы.
Возвращаясь от $u$ к физическому моменту, получаем
\[
 T=\frac{1}{c}[a\sin(x_1+\delta)-k_1x_1-k_2x_2],
 \qquad A_K=\begin{pmatrix}0&1\\-k_1&-(b+k_2)\end{pmatrix}.
\]
При известных $a,c$ компенсация точна. Если вместо них доступны
только $\widehat a,\widehat c$, реализуется
\[
 T=\frac{1}{\widehat c}
 [\widehat a\sin(x_1+\delta)-k_1x_1-k_2x_2].
\]
Для краткости обозначим постоянные
\[
 \Delta_a=\frac{\widehat a c-a\widehat c}{\widehat c},\qquad
 \Delta_c=\frac{c-\widehat c}{\widehat c}.
\]
Прямая подстановка даёт
\[
 \dot{\bm x}=A_K\bm x+B\delta(\bm x),\qquad
 B=\begin{pmatrix}0\\1\end{pmatrix},\qquad
 \delta(\bm x)=\Delta_a\sin(x_1+\delta)
                   -\Delta_c(k_1x_1+k_2x_2).
\]
Буква $\delta$ без аргумента здесь означает заданный угол;
$\delta(\bm x)$ --- векторную ошибку модели из предыдущего подраздела
(в этом примере скалярную).

\textbf{Дополнение для защиты: корректная оценка со слайда 28.}
Из липшицевости синуса следует
$|\sin(x_1+\delta)|\le|x_1|+|\sin\delta|$. Поэтому
\[
 |\delta(\bm x)|\le k\norm{\bm x}+\varepsilon,\qquad
 k=|\Delta_a|+|\Delta_c|\sqrt{k_1^2+k_2^2},\qquad
 \varepsilon=|\Delta_a|\,|\sin\delta|.
\]
На слайде внутри выражения для $\varepsilon$ напечатано
$\sin(x_1+\delta_1)$; такое выражение не является постоянной оценкой
и не даёт указанного там следствия. Правильная постоянная выше
показывает: при $\sin\delta=0$ имеем $\varepsilon=0$.
Для
\[
 P=\begin{pmatrix}p_{11}&p_{12}\\p_{12}&p_{22}\end{pmatrix},\qquad
 PB=\begin{pmatrix}p_{12}\\p_{22}\end{pmatrix}
\]
условие робастной устойчивости принимает вид
\[
 k<\frac{1}{2\sqrt{p_{12}^2+p_{22}^2}}.
\]
Если $\sin\delta\ne0$ и $\Delta_a\ne0$, в заданном положении
остаётся ошибка момента: при $\bm x=0$ возмущение равно
$\Delta_a\sin\delta$. Поэтому утверждать точное удержание заданного
угла без дополнительной компенсации нельзя. Слайд 29 завершает лекцию.

\subsection{Работа 3, задача 1: нулевое собственное значение}

\textit{Дополнение для защиты: подробный разбор задания 1 Work3;
метод --- L4\_2, слайды 3--10.}
Исходная система:
\[
 \dot x_1=-x_1+x_2-x_3,\quad
 \dot x_2=-x_1x_3-x_2+u,\quad
 \dot x_3=-x_1+u,\quad y=x_3.
\]
Здесь $g=(0,1,1)^T$, $h=x_3$, поэтому $L_gh=1$.
Управление появляется сразу: $\rho=1$ во всём $\mathbb R^3$.
Из $\dot y=-x_1+u$ следуют
$\gamma=1$, $\alpha=x_1$ и $u=x_1+v$, дающее $\dot y=v$.

Для внутренней координаты нужно $L_g\varphi=0$, то есть
$\partial\varphi/\partial x_2+\partial\varphi/\partial x_3=0$.
Вход одинаково изменяет $x_2$ и $x_3$, поэтому их разность от него
не зависит; $x_1$ также напрямую от входа не зависит. Отсюда выбор
\[
 \eta_1=x_1,\qquad\eta_2=x_2-x_3,\qquad\xi=x_3.
\]
Это линейная глобальная замена с $\det DT=1$ и обратным
$x_1=\eta_1$, $x_2=\eta_2+\xi$, $x_3=\xi$.
При дифференцировании второй внутренней координаты управление
сокращается:
\[
 \dot\eta_2=\dot x_2-\dot x_3=x_1-x_2-x_1x_3.
\]
Следовательно, нормальная форма исходной системы:
\[
 \dot\eta_1=-\eta_1+\eta_2,\qquad
 \dot\eta_2=\eta_1-\eta_2-(\eta_1+1)\xi,\qquad
 \dot\xi=-\eta_1+u.
\]
На нулевом выходе $\xi=0$ выбираем $u=\eta_1$, то есть $v=0$.
Внутренняя матрица равна
\[
 A_0=\begin{pmatrix}-1&1\\1&-1\end{pmatrix},\qquad
 \lambda_1=0,\quad\lambda_2=-2.
\]
Чтобы понять нулевое значение без формального критерия, введём
$p=\eta_1+\eta_2$, $q=\eta_1-\eta_2$. Тогда $\dot p=0$,
$\dot q=-2q$, и
\[
 \eta_1(t)=\frac{p(0)+q(0)e^{-2t}}{2},\qquad
 \eta_2(t)=\frac{p(0)-q(0)e^{-2t}}{2}.
\]
Разность исчезает, но сумма сохраняется. Поэтому нуль-динамика
устойчива, имеет прямую равновесий $\eta_1=\eta_2$ и не является
асимптотически устойчивой. Вывод: система линеаризуема по входу-выходу,
но не минимально-фазовая. Нулевое собственное значение здесь не
является неопределённым случаем линеаризации: внутренняя модель уже
\emph{точно линейна}, и её решения выписаны явно.

\subsection{Работа 3, задача 2: зачем нужна нелинейная внутренняя координата}

\textit{Дополнение для защиты: подробный разбор задания 2 Work3;
метод --- L4\_1, слайды 10--18, и L4\_2, слайды 6--9.}
Требуется глобально стабилизировать
\[
 \dot x_1=-x_1+x_2,\quad
 \dot x_2=x_1-x_2-x_1x_3+u,\quad
 \dot x_3=x_1+x_1x_2-2x_3.
\]
Выход в условии не задан, поэтому выбираем $y=x_1$.
Тогда $\dot y=-x_1+x_2$,
$\ddot y=2x_1-2x_2-x_1x_3+u$, то есть $\rho=2$ глобально.
Естественные внешние координаты:
\[
 \xi_1=x_1,\qquad\xi_2=-x_1+x_2.
\]
Оставшаяся переменная $x_3$ напрямую не содержит $u$, но её производная
имеет произведение $x_1x_2$. Заметим, что
\[
 \frac{d}{dt}\frac{x_1^2}{2}=x_1(-x_1+x_2)=-x_1^2+x_1x_2.
\]
Вычитание этой производной уничтожает нежелательное произведение:
\[
 \eta=x_3-\frac{x_1^2}{2},\qquad
 \dot\eta=x_1+x_1x_2-2x_3-(-x_1^2+x_1x_2)=\xi_1-2\eta.
\]
Такой выбор усиливает результат по сравнению с одной проверкой
нуль-динамики: внутренняя система становится линейной даже при
ненулевом выходе.
Полная замена имеет
\[
 DT=\begin{pmatrix}-x_1&0&1\\1&0&0\\-1&1&0\end{pmatrix},
 \quad\det DT=1,
 \quad
 x_1=\xi_1,\quad x_2=\xi_1+\xi_2,\quad
 x_3=\eta+\frac{\xi_1^2}{2}.
\]
Это глобальный полиномиальный диффеоморфизм.
Из $\dot\xi_2=-2\xi_2-x_1x_3+u$ выбираем
\[
 u=x_1x_3+2\xi_2+v,\qquad v=-\xi_1-2\xi_2.
\]
Внешний полином равен $(\lambda+1)^2$; после сокращения линейных
слагаемых итоговый закон особенно прост:
\[
 \boxed{u=x_1x_3-x_1}.
\]
Все новые координаты подчиняются одной линейной системе
\[
 \begin{pmatrix}\dot\eta\\\dot\xi_1\\\dot\xi_2\end{pmatrix}
 =\begin{pmatrix}-2&1&0\\0&0&1\\0&-1&-2\end{pmatrix}
   \begin{pmatrix}\eta\\\xi_1\\\xi_2\end{pmatrix}.
\]
Матрица блочно-треугольная с собственными значениями $-2,-1,-1$,
поэтому все координаты стремятся к нулю. Глобальность обеспечивается
глобальной обратимостью замены, отсутствием знаменателей в управлении
и существованием линейного решения при всех $t\ge0$.
Исходное начало глобально асимптотически устойчиво.
На защите важно назвать оба факта: устойчивость внешней цепочки
и устойчивость внутреннего уравнения, а не ограничиться полюсами
$-1,-1$.

\subsection{Работа 3, задача 3: устойчивый выход при расходящемся состоянии}

\textit{Дополнение для защиты: подробный разбор задания 3 Work3;
метод --- L4\_2, слайды 6--10.}
Рассматривается
\[
 \dot x_1=x_2,\quad\dot x_2=-x_1+x_3+u,\quad
 \dot x_3=x_3+x_1^2,\quad y=x_1.
\]
Имеем $L_gh=0$, $L_fh=x_2$, $L_gL_fh=1$,
$L_f^2h=-x_1+x_3$. Значит $\rho=2$ на $\mathbb R^3$,
$\alpha=x_1-x_3$, $\gamma=1$ и
$u=x_1-x_3+v$ даёт $\ddot y=v$.
Внутреннюю координату можно взять непосредственно $\eta=x_3$,
поскольку $g=(0,1,0)^T$ и $L_gx_3=0$.
Замена $(\eta,\xi_1,\xi_2)=(x_3,x_1,x_2)$ --- глобальная перестановка.
Нормальная форма до замены входа:
\[
 \dot\eta=\eta+\xi_1^2,\quad
 \dot\xi_1=\xi_2,\quad
 \dot\xi_2=-\xi_1+\eta+u.
\]
При нулевом выходе остаётся $\dot\eta=\eta$, так что
$\eta(t)=\eta(0)e^t$. Система не минимально-фазовая.

Почему линейная обратная связь $v=-k_1\xi_1-k_2\xi_2$ с
$k_1,k_2>0$ этого не исправляет? При
$\xi_1(0)=\xi_2(0)=0$, $\eta(0)\ne0$ внешняя цепочка остаётся
нулевой, но $x_3(t)=\eta(0)e^t$. Линеаризующий регулятор всё точнее
компенсирует растущее $x_3$ во втором уравнении, сохраняя $y=0$.
Это не стабилизация полного состояния.
Более общая формула
\[
 \eta(t)=e^t\left[\eta(0)+\int_0^t e^{-\tau}\xi_1(\tau)^2\,d\tau\right]
\]
показывает, что даже затухающий внешний процесс способен возбуждать
неустойчивую внутреннюю переменную.

\subsection{Работа 3, задача 4: две ветви координат и вырожденная граница}

\textit{Дополнение для защиты: подробный разбор задания 4 Work3;
метод --- L4\_2, слайды 3--10.}
Исходная система и её поля:
\[
 \dot x_1=x_2,\quad\dot x_2=-x_1+x_3u,\quad
 \dot x_3=-x_3+u,\quad y=x_1,
\]
\[
 f=(x_2,-x_1,-x_3)^T,\qquad g=(0,x_3,1)^T.
\]
Вычисляем $L_gh=0$, $L_fh=x_2$, $L_f^2h=-x_1$,
$L_gL_fh=x_3$. Следовательно,
\[
 \rho=2\quad\text{на }D_0=\{\bm x:x_3\ne0\},\qquad
 u=\frac{x_1+v}{x_3}.
\]
При $x_3=0$ нельзя делить на коэффициент входа. Более того, очередные
чистые производные $L_f^kh$ циклически равны $x_1,x_2,-x_1,-x_2$,
а соответствующие производные вдоль $g$ равны нулю либо $\pm x_3$.
Это не отдельная область с относительной степенью три.

Условие для внутренней координаты имеет вид
$x_3\varphi_{x_2}+\varphi_{x_3}=0$. Вход изменяет $x_2$ со скоростью
$x_3$, а $x_3$ со скоростью $1$, поэтому производная $x_3^2/2$
по входному полю равна $x_3$. Отсюда
\[
 \eta=x_2-\frac{x_3^2}{2},\qquad L_g\eta=x_3-x_3=0,
 \qquad L_f\eta=-x_1+x_3^2.
\]
Для $\xi_1=x_1$, $\xi_2=x_2$ получаем
\[
 T(\bm x)=\begin{pmatrix}x_2-x_3^2/2\\x_1\\x_2\end{pmatrix},\quad
 DT=\begin{pmatrix}0&1&-x_3\\1&0&0\\0&1&0\end{pmatrix},\quad
 \det DT=-x_3.
\]
Преобразование локально обратимо во всех точках $D_0$, но теряет
знак $x_3$. Поэтому глобальную обратную функцию на всём $D_0$
построить нельзя. На каждой ветви $D_\sigma=\{\sigma x_3>0\}$,
$\sigma\in\{-1,1\}$, она существует:
\[
 x_1=\xi_1,\qquad x_2=\xi_2,\qquad
 x_3=\sigma\sqrt{2(\xi_2-\eta)},\qquad \eta<\xi_2.
\]
Обе ветви имеют одно и то же изображение $\{\eta<\xi_2\}$.
Нормальная форма на фиксированной ветви:
\[
 \dot\eta=-\xi_1+2\xi_2-2\eta,\quad
 \dot\xi_1=\xi_2,\quad
 \dot\xi_2=-\xi_1+\sigma\sqrt{2(\xi_2-\eta)}\,u.
\]
При нулевом выходе $\xi_1=\xi_2=0$ необходимо $\eta<0$;
линеаризующая обратная связь с $v=0$ даёт $u=0$ и
\[
 \dot\eta=-2\eta,\qquad\eta(t)=\eta(0)e^{-2t}.
\]
В исходных координатах $\dot x_3=-x_3$,
$x_3(t)=x_3(0)e^{-t}$. Каждая ненулевая ветвь сохраняет знак при
любом конечном времени и приближается к $x_3=0$ лишь при $t\to\infty$.
Продолженная на эту точку нуль-динамика экспоненциально устойчива.
Однако сама точка $\eta=0$, $\bm\xi=0$ лежит на границе области
координат, а $u=(x_1+v)/x_3$ в ней не определено.
Поэтому устойчивость продолженной нуль-динамики не доказывает
глобальную стабилизацию этим законом около исходного начала и
не позволяет без оговорок применять там стандартное определение
минимально-фазовости с регулярным $T$.

\subsection{Работа 3, задача 5: полная цепочка и восстановление управления}

\textit{Дополнение для защиты: подробный разбор задания 5 Work3;
метод --- L4\_1, слайды 10--18, и L4\_2, слайды 3--8.}
Здесь обозначения $\bm z,\alpha,\beta$ сохранены непосредственно
из условия задания:
\[
 \dot x_1=x_2,\quad\dot x_2=x_3+x_1^2,\quad
 \dot x_3=-x_1x_2+(1+x_1^2)u,\quad y=x_1.
\]
Поля равны $f=(x_2,x_3+x_1^2,-x_1x_2)^T$,
$g=(0,0,1+x_1^2)^T$. Последовательно получаем
\[
 L_gh=0,\quad L_fh=x_2,\quad L_gL_fh=0,\quad
 L_f^2h=x_3+x_1^2,
\]
\[
 L_f^3h=-x_1x_2+2x_1x_2=x_1x_2,\qquad
 L_gL_f^2h=1+x_1^2>0.
\]
Поэтому $\rho=3=n$ на всей $\mathbb R^3$. Все состояния входят
во внешнюю цепочку, внутренних координат и нуль-динамики положительной
размерности нет. Выбор
\[
 z_1=x_1,\qquad z_2=x_2,\qquad z_3=x_3+x_1^2
\]
даёт $\det DT=1$ и глобальное обратное
$x_1=z_1$, $x_2=z_2$, $x_3=z_3-z_1^2$. Уравнения:
\[
 \dot z_1=z_2,\qquad\dot z_2=z_3,\qquad
 \dot z_3=x_1x_2+(1+x_1^2)u.
\]
Отсюда
\[
 u=\alpha(\bm x)+\beta(\bm x)v,
 \qquad\alpha=-\frac{x_1x_2}{1+x_1^2},\qquad
 \beta=\frac{1}{1+x_1^2}.
\]
В обозначениях лекции $\beta=\gamma^{-1}$: это коэффициент перед
новым входом, а не усиление скользящего регулятора из L6.
После компенсации $\dot{\bm z}=A\bm z+Bv$ с
\[
 A=\begin{pmatrix}0&1&0\\0&0&1\\0&0&0\end{pmatrix},\qquad
 B=\begin{pmatrix}0\\0\\1\end{pmatrix}.
\]
Для $v=-K\bm z$, $K=(k_1,k_2,k_3)$,
\[
 A-BK=\begin{pmatrix}0&1&0\\0&0&1\\-k_1&-k_2&-k_3\end{pmatrix},
 \qquad p(\lambda)=\lambda^3+k_3\lambda^2+k_2\lambda+k_1.
\]
Требуемые полюса задают
\[
 (\lambda+2)^2(\lambda+3)=\lambda^3+7\lambda^2+16\lambda+12,
 \qquad K=\begin{pmatrix}12&16&7\end{pmatrix}.
\]
Сначала записываем $v=-12z_1-16z_2-7z_3$, затем подставляем
\emph{все} выражения $z_i(\bm x)$:
\[
 \boxed{u(\bm x)=
 \frac{-x_1x_2-12x_1-16x_2-7x_3-7x_1^2}{1+x_1^2}}.
\]
Слагаемое $-7x_1^2$ возникает из $-7z_3$; его легко потерять,
если слишком рано возвращаться к исходным переменным.
Знаменатель никогда не обращается в нуль, обратное преобразование
глобально, а линейная матрица гурвицева. Поэтому закон определён
глобально и даёт глобальную асимптотическую устойчивость исходного
начала. Кратный полюс $-2$ допускает множитель $t e^{-2t}$ в решении,
который также стремится к нулю.

\subsection{Как объяснить ход решения на защите}

\textit{Дополнение для защиты: объединение алгоритмов L4\_1 и L4\_2.}
Для конкретной системы последовательно отвечаем на шесть вопросов:
каков выход; когда управление впервые появляется в его производных;
где последний коэффициент отличен от нуля; какие внутренние
координаты удовлетворяют $L_g\varphi_i=0$; обратима ли полная замена;
что происходит с внутренними состояниями после линеаризации выхода.
Лишь после этого выбираем $K$ и возвращаем управление к исходным
координатам. В задаче 1 препятствие --- нулевая мода внутренней
системы, в задаче 3 --- её неустойчивая мода, в задаче 4 ---
вырождение преобразования и обратной связи. В задачах 2 и 5
глобальный результат обоснован отдельной проверкой всех новых
координат и глобального обратного преобразования.
